%bismi ALLAH ARRAHMAN ARRAHIM

































































\documentclass[11pt,a4paper,oneside]{book}

% ============================================================
% THÉORIE DES EMPREINTES
% Extension : cerveaux de dataset, maillage global,
% écosystèmes causaux, contexte maillé et recherche de réponse
%
% Statut :
%   - socle observationnel hérité de la Théorie des Empreintes ;
%   - nouvelles définitions propres à la présente extension ;
%   - propositions démontrées sous hypothèses explicites ;
%   - postulats opérationnels et conditions de validation.
% ============================================================

\usepackage[utf8]{inputenc}
\usepackage[T1]{fontenc}
\usepackage[french]{babel}
\usepackage{lmodern}
\usepackage{microtype}
\usepackage[a4paper,margin=2.6cm]{geometry}
\usepackage{amsmath,amssymb,amsthm,mathtools}
\usepackage{bm}
\usepackage{enumitem}
\usepackage{booktabs}
\usepackage{longtable}
\usepackage{array}
\usepackage{xcolor}
\usepackage{hyperref}
\usepackage{csquotes}

\hypersetup{
  colorlinks=true,
  linkcolor=blue!45!black,
  citecolor=blue!45!black,
  urlcolor=blue!45!black,
  pdftitle={Cerveaux contextuels, maillage global et recherche causale
  dans la Théorie des Empreintes},
  pdfauthor={Extension théorique de la Théorie des Empreintes}
}

\setcounter{secnumdepth}{3}
\setcounter{tocdepth}{3}
\setlength{\parskip}{0.35em}
\setlength{\parindent}{1.2em}

\newtheorem{definition}{Définition}[chapter]
\newtheorem{postulat}[definition]{Postulat}
\newtheorem{hypothese}[definition]{Hypothèse}
\newtheorem{proposition}[definition]{Proposition}
\newtheorem{theoreme}[definition]{Théorème}
\newtheorem{lemme}[definition]{Lemme}
\newtheorem{corollaire}[definition]{Corollaire}
\theoremstyle{remark}
\newtheorem{remarque}[definition]{Remarque}
\newtheorem{limite}[definition]{Limite}
\newtheorem{exemple}[definition]{Exemple}

\DeclareMathOperator{\Imat}{Im}
\DeclareMathOperator{\Dom}{Dom}
\DeclareMathOperator{\Supp}{Supp}
\DeclareMathOperator{\Path}{Path}
\DeclareMathOperator{\Reach}{Reach}
\DeclareMathOperator{\cost}{cost}
\DeclareMathOperator{\Conf}{Conf}
\DeclareMathOperator{\Align}{Align}
\DeclareMathOperator{\Compat}{Compat}
\DeclareMathOperator{\Cand}{Cand}
\DeclareMathOperator{\Score}{Score}
\DeclareMathOperator{\Cl}{Cl}
\DeclareMathOperator{\Fib}{Fib}
\DeclareMathOperator{\Tok}{Tok}
\DeclareMathOperator{\Occ}{Occ}
\DeclareMathOperator{\Eco}{Eco}
\DeclareMathOperator{\Brain}{Brain}
\DeclareMathOperator{\Ctx}{Ctx}
\DeclareMathOperator{\Ans}{Ans}
\DeclareMathOperator{\Obs}{Obs}
\DeclareMathOperator{\Cert}{Cert}
\DeclareMathOperator{\Prov}{Prov}
\DeclareMathOperator{\Hard}{Hard}
\DeclareMathOperator{\Soft}{Soft}

\newcommand{\R}{\mathbb{R}}
\newcommand{\N}{\mathbb{N}}
\newcommand{\Rp}{\mathbb{R}_{\geq 0}}
\newcommand{\Bool}{\{0,1\}}
\newcommand{\Nat}{\mathsf{Nat}}
\newcommand{\Anti}{\mathsf{Anti}}
\newcommand{\Neutre}{\mathsf{Neutre}}
\newcommand{\Unknown}{\mathsf{Unknown}}
\newcommand{\Global}{\mathcal{G}}
\newcommand{\Dataset}{\mathcal{D}}
\newcommand{\Prompts}{\mathcal{P}}
\newcommand{\Answers}{\mathcal{R}}
\newcommand{\Tokens}{\mathcal{V}}
\newcommand{\Edges}{\mathcal{E}}
\newcommand{\Rel}{\mathcal{L}}
\newcommand{\Causal}{\mathsf{causal}}
\newcommand{\Anticausal}{\mathsf{anticausal}}
\newcommand{\ObsRel}{\sim_{\mathrm{obs}}}
\newcommand{\TokRel}{\sim_{\mathrm{tok}}}
\newcommand{\eps}{\varepsilon}

\title{
  \textbf{Théorie des Empreintes Informationnelles}\\[0.4em]
  \large Cerveaux de dataset, maillage global,\\
  écosystèmes causaux et contexte maillé\\[0.4em]
  \normalsize Formalisation d'une architecture de recherche\\
  et de construction de réponses contextuellement alignées
}
\author{Extension théorique du socle des empreintes}
\date{\today}

\begin{document}

\maketitle
\tableofcontents

% ============================================================
\part{Statut, portée et méthode}
% ============================================================

\chapter{Objet et statut de la présente extension}

\section{Objet}

La présente extension formalise un mécanisme de construction et
d'exploitation de structures relationnelles dans un système fondé sur
la Théorie des Empreintes Informationnelles.

Elle traite conjointement :

\begin{enumerate}[label=\arabic*.]
  \item la constitution d'un cerveau local pour chaque exemple
        prompt--réponse d'un dataset ;
  \item la construction d'un maillage global par agrégation des
        cerveaux locaux et fusion des identités lexicales communes ;
  \item la conservation des écosystèmes relationnels associés aux
        tokens, notamment des relations dites causales naturelles
        et anti-causales ;
  \item la construction d'un maillage contextuel propre à chaque
        nouveau prompt ;
  \item la recherche, dans le maillage global, d'une trajectoire
        compatible avec le maillage contextuel ;
  \item l'alignement des éléments du prompt avec des tokens, des
        ensembles de tokens ou des trajectoires intermédiaires ;
  \item l'orientation du processus de recherche par une trajectoire
        causale globale, notée notamment
        \(c_0\leadsto c_x\) ;
  \item la discrimination des réponses candidates par la
        combinaison de l'alignement contextuel, de la structure
        relationnelle globale et des contraintes causales.
\end{enumerate}

L'objectif est de définir une architecture dans laquelle une réponse
n'est pas retenue uniquement parce qu'elle contient des tokens
associés à la question. Elle doit pouvoir être rattachée à une
trajectoire du maillage global qui satisfait les contraintes
contextuelles et causales du protocole.

\section{Sources théoriques et règle de fidélité}

Le socle mathématique utilisé est celui de la Théorie des Empreintes
Informationnelles, notamment ses constructions relatives :

\begin{itemize}
  \item aux objets observables, aux contextes et aux sondes ;
  \item à l'empreinte informationnelle ;
  \item à l'indiscernabilité, aux fibres et aux quotients ;
  \item à la factorisation des observations ;
  \item aux échelles et aux seuils de résolution ;
  \item aux trajectoires, à la généralisation observationnelle,
        à la recherche et à la séparation des niveaux de traitement ;
  \item à la distinction entre prédiction conditionnelle et causalité
        établie.
\end{itemize}

Le présent texte ajoute un formalisme de graphes lexicaux et de
recherche contextuelle. Ces constructions ne sont pas présentées
comme si elles figuraient déjà intégralement dans le socle
observationnel.

\section{Statuts logiques}

Les statuts suivants sont distingués.

\begin{definition}[Statuts d'un énoncé]
Un énoncé de la présente extension est classé comme suit :
\begin{description}[leftmargin=3.2cm,style=nextline]
  \item[Définition] Fixe la signification d'un objet ou d'une opération.
  \item[Postulat] Introduit une règle constitutive de l'architecture.
  \item[Hypothèse] Énonce une condition requise pour un résultat.
  \item[Proposition] Énonce une propriété démontrable à partir des
        définitions et hypothèses.
  \item[Théorème] Énonce un résultat démontré dans le domaine déclaré.
  \item[Convention] Fixe un choix de représentation ou d'implémentation.
  \item[Résultat empirique] Dépend d'une observation ou d'une expérience.
  \item[Conjecture] Désigne une propriété qui n'est pas démontrée.
\end{description}
\end{definition}

\begin{remarque}
Une définition ne prouve pas l'efficacité de l'architecture qu'elle
décrit. Une proposition démontrée sur un graphe fini ne démontre pas
automatiquement une propriété universelle de tous les systèmes de
langage. Les résultats sont toujours rapportés à leur domaine
d'application.
\end{remarque}


% ============================================================
\chapter{Raccordement au socle de la Théorie des Empreintes}
% ============================================================

\section{Objets observables, contextes et sondes}

Soient :

\begin{itemize}
  \item \(\mathcal{X}\), un ensemble d'objets observables ;
  \item \(\mathcal{P}\), un ensemble de contextes d'observation ;
  \item \(\mathcal{F}\), une famille de sondes ;
  \item \(\Phi\), une application d'observation.
\end{itemize}

On considère une application de la forme

\[
\Phi:\mathcal{P}\times\mathcal{X}\times\mathcal{F}
\longrightarrow \mathcal{Y},
\]

où \(\mathcal{Y}\) est l'espace des réponses des sondes.

Pour un contexte \(P\), un objet \(x\) et une famille finie de sondes
\(\mathcal{F}_K=\{f_1,\ldots,f_K\}\), l'empreinte est

\[
I_\Phi^P(x)
=
\bigl(
\Phi(P,x,f_1),\ldots,\Phi(P,x,f_K)
\bigr).
\]

Plus généralement, pour une famille de sondes \(\mathcal{G}\),

\[
I_{\Phi,\mathcal{G}}^P(x)
=
\bigl(f\mapsto \Phi(P,x,f)\bigr)_{f\in\mathcal{G}}.
\]

\begin{definition}[Indiscernabilité observationnelle]
Pour un contexte \(P\) et une famille de sondes \(\mathcal{G}\),
on définit

\[
x\ObsRel_{P,\mathcal{G}} y
\quad\Longleftrightarrow\quad
I_{\Phi,\mathcal{G}}^P(x)
=
I_{\Phi,\mathcal{G}}^P(y).
\]

La classe d'équivalence de \(x\) est sa fibre observationnelle :

\[
[x]_{P,\mathcal{G}}
=
\{y\in\mathcal{X}:y\ObsRel_{P,\mathcal{G}}x\}.
\]
\end{definition}

La représentation observable identifie ainsi les objets que les
sondes disponibles ne savent pas distinguer. La théorie ne suppose
pas qu'une représentation encode automatiquement toutes les
propriétés d'un objet.

\section{Factorisation et portée de l'information}

\begin{proposition}[Factorisation par l'empreinte]
Soit \(I:\mathcal{X}\to\mathcal{Z}\) une empreinte et
\(A:\mathcal{X}\to\mathcal{Y}\) une application.

Si \(A\) est constante sur les fibres de \(I\), alors il existe une
application unique
\(\overline{A}:I(\mathcal{X})\to\mathcal{Y}\) telle que

\[
A=\overline{A}\circ I.
\]

\end{proposition}

\begin{proof}
Pour \(z\in I(\mathcal{X})\), choisissons \(x\) tel que \(I(x)=z\)
et posons \(\overline{A}(z)=A(x)\). Si \(x'\) est un autre représentant
de \(z\), alors \(I(x)=I(x')\). La constance de \(A\) sur les fibres
donne \(A(x)=A(x')\), donc \(\overline{A}\) est bien définie.
L'unicité découle de la surjectivité de \(I\) sur son image.
\end{proof}

\begin{remarque}[Conséquence pour les cerveaux de langage]
Le cerveau lexical d'un exemple, le maillage global et le cerveau
contextuel sont des structures de représentation. Leur existence ne
prouve pas, à elle seule, qu'ils préservent toute distinction
sémantique pertinente.

Pour établir une propriété de préservation, il faut préciser les
sondes, les relations observées, les critères d'équivalence et la
portée de la factorisation.
\end{remarque}

\section{Causalité et prédiction : distinction héritée}

Une fréquence conditionnelle, une association observée ou une
transition fréquente ne devient pas une causalité au seul motif
qu'on la nomme ainsi.

Dans la présente extension, les termes \emph{causal naturel} et
\emph{anti-causal} désignent des catégories relationnelles internes
au modèle. Leur usage est opérationnel : les arêtes reçoivent un type
causal par une procédure de construction ou d'annotation explicitée.

Ces catégories ne sont pas assimilées automatiquement à une causalité
physique ou à une causalité interventionnelle. Une telle interprétation
demanderait des hypothèses ou des mécanismes supplémentaires.


% ============================================================
\part{Cerveaux locaux et construction du dataset}
% ============================================================

\chapter{Représentation des exemples prompt--réponse}

\section{Token, occurrence et identité lexicale}

\begin{definition}[Token canonique]
Un token canonique est une unité lexicale identifiée selon une règle
de tokenisation et de normalisation déclarée.

On note \(\Tokens\) l'ensemble des identités lexicales canoniques.

La normalisation peut prendre en compte, selon le protocole :
la casse, la ponctuation, les espaces, les variantes Unicode, les
marqueurs de sous-mots et les conventions propres au tokenizer.
\end{definition}

\begin{definition}[Occurrence de token]
Une occurrence est un événement localisé dans un exemple. Elle est
représentée par

\[
o=(z,s,i,v),
\]

où \(z\) identifie l'exemple du dataset, \(s\) le segment
(\(\mathsf{prompt}\) ou \(\mathsf{réponse}\)), \(i\) la position
dans le segment et \(v\in\Tokens\) l'identité canonique du token.

L'application d'identité lexicale est

\[
\pi_{\mathrm{tok}}(o)=v.
\]
\end{definition}

\begin{remarque}
Deux occurrences peuvent partager la même identité lexicale sans
être confondues comme occurrences. La fusion porte sur l'identité
canonique ; la position, le rôle local, la provenance et les relations
propres à chaque occurrence doivent rester reconstructibles.
\end{remarque}

\section{Exemples du dataset}

\begin{definition}[Exemple prompt--réponse]
Un exemple de dataset est un tuple

\[
z=(P_z,R_z,m_z),
\]

où :
\begin{itemize}
  \item \(P_z\) est le prompt ;
  \item \(R_z\) est la réponse associée ;
  \item \(m_z\) est une métadonnée de provenance et de statut.
\end{itemize}

Les métadonnées peuvent notamment identifier la source, la version,
la date, le statut de validation, la qualité de l'annotation et les
conditions d'admission de l'exemple dans le maillage global.
\end{definition}

Le dataset est un ensemble ou un multiensemble fini d'exemples :

\[
\Dataset=\{z_1,\ldots,z_N\}.
\]

La convention de multiensemble est utile lorsque plusieurs occurrences
d'un même couple prompt--réponse doivent conserver leur multiplicité
empirique.

\section{Le cerveau local d'un exemple}

\begin{definition}[Cerveau local d'un exemple]
Le cerveau local associé à \(z\), noté \(\Brain_z\), est une structure

\[
\Brain_z=
\bigl(
V_z,E_z,\tau_z,\ell_z,w_z,\operatorname{prov}_z
\bigr),
\]

où :

\begin{itemize}
  \item \(V_z\) est l'ensemble fini des nœuds locaux ;
  \item \(E_z\subseteq V_z\times V_z\) est l'ensemble des arêtes
        orientées ;
  \item \(\tau_z:V_z\to\mathcal{T}\) donne le type des nœuds ;
  \item \(\ell_z:E_z\to\mathcal{L}\) donne le type des relations ;
  \item \(w_z:E_z\to\R_{\geq 0}\) donne un poids de soutien
        lorsqu'un poids est défini ;
  \item \(\operatorname{prov}_z\) conserve la provenance des nœuds,
        arêtes et annotations.
\end{itemize}
\end{definition}

Les types de nœuds peuvent inclure :

\[
\mathcal{T}=
\{
\mathsf{token},
\mathsf{occurrence},
\mathsf{segment},
\mathsf{slot},
\mathsf{état},
\mathsf{trajectoire}
\}.
\]

Les types d'arêtes peuvent inclure :

\[
\mathcal{L}=
\{
\mathsf{ordre},
\mathsf{contexte},
\mathsf{alignement},
\mathsf{causal},
\mathsf{anticausal},
\mathsf{contraste},
\mathsf{provenance}
\}.
\]

Cette liste est une signature minimale extensible. L'ajout d'un type
de relation exige une définition de sa signification et de ses
conditions d'utilisation.

\section{Construction du cerveau local}

Soient les tokenisations

\[
P_z=(p_1,\ldots,p_m),\qquad
R_z=(r_1,\ldots,r_n).
\]

La construction de \(\Brain_z\) s'effectue selon les opérations
suivantes.

\begin{enumerate}[label=\textbf{Étape \arabic*.}]
  \item \textbf{Tokenisation.}
        Appliquer un tokenizer déclaré au prompt et à la réponse.

  \item \textbf{Création des occurrences.}
        Créer un nœud d'occurrence distinct pour chaque position
        de chaque segment.

  \item \textbf{Enregistrement des identités.}
        Associer chaque occurrence à son token canonique sans
        fusionner les occurrences.

  \item \textbf{Relations internes.}
        Enregistrer les relations observables et les relations
        définies par le protocole : ordre local, relations
        contextuelles, alignements et autres relations autorisées.

  \item \textbf{Trajectoire de réponse.}
        Lorsque l'exemple est admis comme exemple prompt--réponse,
        enregistrer la relation entre le contexte du prompt et
        la trajectoire conduisant à la réponse associée.

  \item \textbf{Écosystèmes.}
        Construire les relations typées associées aux tokens à partir
        des observations ou annotations disponibles.

  \item \textbf{Provenance et statut.}
        Conserver l'identité de l'exemple, le statut de validation,
        les annotations et la méthode de construction de chaque
        relation.
\end{enumerate}

\begin{remarque}[Absence d'inférence non justifiée]
La présence de deux tokens dans un même exemple ne suffit pas à
créer automatiquement une relation causale entre eux.

Une relation doit être produite par une règle de construction,
une annotation, une observation ou un mécanisme déclaré. Si le
protocole ne permet pas de décider son type, le type reste
\(\Unknown\), plutôt que d'être arbitrairement déclaré causal.
\end{remarque}

\section{Brain de prompt, brain de réponse et brain d'exemple}

On distingue trois sous-structures :

\[
\Brain_z^P,\qquad
\Brain_z^R,\qquad
\Brain_z^{PR}.
\]

\begin{itemize}
  \item \(\Brain_z^P\) représente les occurrences et relations
        internes au prompt ;
  \item \(\Brain_z^R\) représente les occurrences et relations
        internes à la réponse ;
  \item \(\Brain_z^{PR}\) représente les relations documentées
        entre le prompt et sa réponse.
\end{itemize}

Le cerveau complet peut être obtenu par recollement de ces
sous-structures sur leurs interfaces déclarées :

\[
\Brain_z=
\Brain_z^P\cup
\Brain_z^R\cup
\Brain_z^{PR}.
\]

Cette union conserve les identifiants d'occurrences, les types
d'arêtes et la provenance. Elle ne fusionne pas deux occurrences
simplement parce que leurs formes lexicales sont identiques.

\begin{postulat}[Traçabilité des relations prompt--réponse]
Toute relation entre un prompt et une réponse dans un cerveau local
doit être traçable jusqu'à l'exemple source et au protocole qui
l'a produite.

L'admission d'un exemple dans le dataset n'implique pas, à elle
seule, que sa réponse soit factuellement vraie. Le statut de validation
de l'exemple doit rester disponible.
\end{postulat}


% ============================================================
\chapter{Écosystèmes des tokens}
% ============================================================

\section{Définition de l'écosystème}

\begin{definition}[Écosystème d'un token]
Soit \(G\) un maillage global et \(v\in\Tokens\) un token canonique.
L'écosystème de \(v\), noté \(\Eco_G(v)\), est la structure relationnelle
incidente à \(v\), conservant les types, directions, poids et
provenances des relations.

Pour une version strictement locale :

\[
\Eco_G^{(1)}(v)
=
\{(u,r,w): (v,u)\in E_G
\text{ ou }(u,v)\in E_G\},
\]

où \(r\) est le type de relation et \(w\) son poids éventuel.

Pour une version étendue de profondeur \(k\), on définit
\(\Eco_G^{(k)}(v)\) comme le sous-graphe induit par les nœuds
accessibles à distance relationnelle au plus \(k\), selon les types
d'arêtes autorisés.
\end{definition}

L'écosystème n'est donc pas seulement une liste de tokens associés.
Il conserve la structure typée des relations, notamment leur
orientation et leur provenance.

\section{Causalité naturelle et anti-causalité}

\begin{definition}[Relation causale naturelle interne]
Une arête
\[
u\xrightarrow{\Nat}v
\]
est une relation classée comme causale naturelle par le protocole
d'annotation ou de construction du maillage.

Le type \(\Nat\) est une étiquette relationnelle interne au modèle.
Son attribution doit être traçable. Elle ne signifie pas, sans
hypothèse supplémentaire, qu'une causalité physique a été démontrée.
\end{definition}

\begin{definition}[Relation anti-causale interne]
Une arête
\[
u\xrightarrow{\Anti}v
\]
est une relation classée dans la catégorie anti-causale de la théorie.

Cette catégorie est distincte de la relation causale naturelle.
Elle n'est pas définie comme le simple renversement de la flèche
\(u\to v\), ni comme une négation logique automatique de cette
relation.
\end{definition}

\begin{remarque}[Orientation et type d'arête]
Il faut distinguer :

\begin{enumerate}
  \item l'orientation graphique d'une arête ;
  \item le type de relation porté par l'arête ;
  \item l'ordre causal d'une trajectoire complète.
\end{enumerate}

Ainsi, une arête orientée \(u\to v\) de type anti-causal reste
orientée de \(u\) vers \(v\) dans le graphe. Son type n'est pas
déduit de la seule orientation.
\end{remarque}

\section{Exemple lexical de référence}

Considérons les contextes :

\begin{quote}
  \texttt{je vais bien merci}\\
  \texttt{c'est bien de le dire}\\
  \texttt{oui bien que}
\end{quote}

Pour le token \texttt{bien}, la structure illustrative fournie
par la théorie est :

\[
\Eco(\texttt{bien})
=
\Eco_{\Nat}(\texttt{bien})
\cup
\Eco_{\Anti}(\texttt{bien}),
\]

avec :

\[
\Eco_{\Nat}(\texttt{bien})
=
\{
\texttt{merci},
\texttt{de},
\texttt{que}
\},
\]

et

\[
\Eco_{\Anti}(\texttt{bien})
=
\{
\texttt{c'est},
\texttt{oui},
\texttt{vais}
\}.
\]

Dans la notation orientée :

\[
\begin{aligned}
\texttt{bien}&\xrightarrow{\Nat}\texttt{merci},\\
\texttt{bien}&\xrightarrow{\Nat}\texttt{de},\\
\texttt{bien}&\xrightarrow{\Nat}\texttt{que},
\end{aligned}
\qquad
\begin{aligned}
\texttt{bien}&\xrightarrow{\Anti}\texttt{c'est},\\
\texttt{bien}&\xrightarrow{\Anti}\texttt{oui},\\
\texttt{bien}&\xrightarrow{\Anti}\texttt{vais}.
\end{aligned}
\]

\begin{remarque}
Cet exemple fixe les catégories et les liens de l'exemple théorique.
Il ne constitue pas une preuve que ces catégories sont universelles
pour tous les usages du mot \texttt{bien}. Leur généralisation exige
un protocole d'annotation et de validation sur les contextes
effectivement étudiés.
\end{remarque}

\section{Rôle fonctionnel de l'anti-causalité}

Pour rendre l'anti-causalité opératoire, il faut spécifier comment
elle influence la recherche. La présente extension distingue
deux fonctions possibles, qui peuvent coexister.

\begin{definition}[Anti-causalité bloquante]
Une relation anti-causale est bloquante dans un contexte donné
lorsqu'une règle de compatibilité certifiée interdit un parcours
qui la viole.

On note
\[
e\in\Hard
\]
le statut d'une telle contrainte.
\end{definition}

\begin{definition}[Anti-causalité réorientante]
Une relation anti-causale est réorientante lorsqu'elle constitue
un signal de contraste ou d'incompatibilité préférentielle, utilisé
pour explorer d'autres branches ou réduire la priorité d'une
trajectoire sans l'interdire absolument.

On note
\[
e\in\Soft
\]
le statut de ce type de relation.
\end{definition}

\begin{postulat}[Traitement conditionnel des relations anti-causales]
Le traitement d'une relation anti-causale dépend de son statut :

\begin{itemize}
  \item une anti-causalité bloquante élimine les trajectoires
        incompatibles avec la contrainte certifiée ;
  \item une anti-causalité réorientante modifie la priorité de
        recherche et favorise l'exploration d'autres trajectoires ;
  \item une relation dont le statut n'est pas établi ne peut pas
        être transformée silencieusement en interdiction absolue.
\end{itemize}

Le protocole doit conserver la provenance de la décision et son
statut de certification.
\end{postulat}

Cette séparation permet à l'anti-causalité d'intervenir à la fois
dans l'exclusion de parcours incompatibles et dans la réorientation
de la recherche, sans confondre les deux opérations.


% ============================================================
\chapter{Fusion des cerveaux locaux dans le maillage global}
% ============================================================

\section{Identité lexicale et fusion}

La fusion des tokens communs constitue la règle d'agrégation du
maillage global. Elle porte sur les identités lexicales canoniques,
et non sur la totalité des occurrences ou de leurs contextes.

\begin{definition}[Équivalence lexicale de fusion]
Sur les nœuds qui représentent des identités lexicales canoniques,
on définit :

\[
u\TokRel v
\quad\Longleftrightarrow\quad
\pi_{\mathrm{tok}}(u)=\pi_{\mathrm{tok}}(v).
\]

La fusion identifie les nœuds de même identité lexicale selon cette
relation d'équivalence.
\end{definition}

\begin{definition}[Maillage global]
Pour un dataset \(\Dataset\), le maillage global \(G_{\Dataset}\)
est la structure relationnelle agrégée des cerveaux locaux :

\[
G_{\Dataset}
=
\operatorname{Aggregate}
\bigl(
\Brain_{z_1},\ldots,\Brain_{z_N}
\bigr).
\]

Il est représenté par :

\[
G_{\Dataset}
=
(V_G,E_G,\tau_G,\ell_G,w_G,\operatorname{Prov}_G).
\]

Les nœuds lexicaux canoniques sont fusionnés par identité. Les
occurrences locales, les contextes d'origine, les types de relations
et les provenances demeurent accessibles.
\end{definition}

\section{Agrégation des relations}

Soit \(e=(u,r,v)\) une relation orientée de type \(r\), observée dans
un ou plusieurs cerveaux locaux.

On peut définir un poids de soutien empirique :

\[
w_G(u,r,v)
=
\sum_{z\in\Dataset}
\omega_z\,n_z(u,r,v),
\]

où :

\begin{itemize}
  \item \(n_z(u,r,v)\) est le nombre de contributions admissibles
        de l'exemple \(z\) à la relation considérée ;
  \item \(\omega_z\geq 0\) est le poids déclaré de l'exemple.
\end{itemize}

Le poids peut être un comptage, une fréquence, une confiance
d'annotation ou une autre quantité définie par le protocole.

Si les types relationnels diffèrent, ils doivent rester distincts.
On ne doit pas fusionner une relation causale naturelle et une
relation anti-causale en une relation non typée unique.

\begin{definition}[Provenance d'une arête globale]
La provenance d'une arête globale est l'ensemble des contributions
locales admissibles qui la soutiennent :

\[
\operatorname{Prov}_G(e)
=
\{(z,e_z,\operatorname{statut}_z):
e_z\text{ contribue à }e\}.
\]

Une provenance peut inclure la source, le statut de validation,
la méthode d'annotation et le contexte d'observation.
\end{definition}

\section{Fusion sans effacement des écosystèmes}

\begin{proposition}[Conservation de la provenance sous fusion]
Supposons que l'opération d'agrégation conserve pour chaque arête
globale la liste des contributions locales, leurs types et leurs
provenances. Alors la fusion des identités lexicales ne rend pas
irrécupérables les contributions relationnelles locales.

\end{proposition}

\begin{proof}
Chaque contribution locale est conservée comme élément de
\(\operatorname{Prov}_G(e)\). L'identification des nœuds lexicaux
ne supprime pas cette liste. Les contributions peuvent donc être
retrouvées à partir de l'arête globale et de sa provenance.
\end{proof}

\begin{remarque}[Deux niveaux à ne pas confondre]
La fusion lexicale n'est pas la fusion des occurrences. Si le token
\texttt{bien} apparaît dans cent exemples, le maillage global peut
posséder une identité lexicale canonique pour \texttt{bien}, mais
doit conserver les cent occurrences et leurs contextes comme
contributions distinctes.

Cette règle permet de partager l'écosystème du token sans effacer
la diversité des contextes qui l'ont constitué.
\end{remarque}

\section{Propriétés minimales de l'agrégation}

\begin{postulat}[Fidélité de l'agrégation]
L'opération d'agrégation doit satisfaire les conditions suivantes :

\begin{enumerate}
  \item les identités lexicales communes sont fusionnées selon
        la règle de normalisation déclarée ;
  \item les types de relations ne sont pas confondus ;
  \item la direction des arêtes est conservée ;
  \item les contributions locales restent traçables ;
  \item les relations inconnues ne sont pas promues silencieusement
        en relations causales ;
  \item les statuts de validation des exemples restent accessibles ;
  \item l'ajout d'un exemple ne doit pas modifier rétroactivement
        la provenance des exemples précédents.
\end{enumerate}
\end{postulat}

\begin{remarque}
La présente extension utilise une agrégation append-only comme
convention de traçabilité : les nouvelles contributions s'ajoutent
au maillage, tandis que toute correction ou invalidation est
enregistrée comme un événement distinct. Une vue courante peut
exclure certaines contributions, mais la provenance doit permettre
de reconstruire la vue et la politique utilisées.
\end{remarque}


% ============================================================
\part{Maillage contextuel du prompt}
% ============================================================

\chapter{Construction du contexte maillé}

\section{Séparation entre maillage global et maillage du prompt}

Deux structures distinctes sont désormais définies :

\[
G_{\Dataset}
\qquad\text{et}\qquad
C_P.
\]

Le maillage global \(G_{\Dataset}\) est construit à partir du dataset
et de la fusion des identités lexicales communes.

Le maillage contextuel \(C_P\) est construit pour le prompt courant
\(P\). Il représente les relations contextuelles nécessaires pour
orienter la recherche d'une réponse.

\begin{postulat}[Non-fusion du cerveau contextuel avec le maillage global]
Le cerveau contextuel d'un prompt est une structure de requête
distincte. Il n'est pas intégré structurellement au maillage global
du seul fait qu'il est construit.

La correspondance entre ses tokens et les identités lexicales du
maillage global sert à rechercher des alignements et des trajectoires.
Elle ne modifie pas les nœuds, les arêtes ou les poids du maillage
global.
\end{postulat}

Il faut donc distinguer :

\begin{enumerate}
  \item la fusion structurelle des cerveaux du dataset, utilisée
        lors de la construction du maillage global ;
  \item la mise en correspondance du prompt courant avec le maillage
        global, utilisée lors de l'inférence.
\end{enumerate}

La seconde opération n'est pas une nouvelle fusion globale.

\section{Tokenisation du prompt}

Soit un prompt \(P\) tokenisé sous la forme :

\[
P=(x_1,x_2,\ldots,x_n),
\qquad x_i\in\Tokens.
\]

Chaque position \(i\) possède une occurrence propre, même lorsque
plusieurs positions contiennent le même token canonique.

On note :

\[
O_P=\{o_1,\ldots,o_n\}
\]

l'ensemble des occurrences du prompt.

\section{Variables contextuelles \(c_i\)}

\begin{definition}[Variable contextuelle]
Une variable contextuelle \(c_i\) est un emplacement de recherche
dans le cerveau contextuel. Lors de la recherche dans le maillage
global, elle peut être instanciée par un objet admissible parmi :

\begin{enumerate}
  \item un token canonique ;
  \item un ensemble fini de tokens ;
  \item une trajectoire finie de tokens ;
  \item un sous-graphe fini ;
  \item une structure relationnelle finie autorisée par le protocole.
\end{enumerate}

L'ensemble des valeurs admissibles est noté \(\mathcal{A}(G)\).
\end{definition}

Ainsi, un \(c_i\) ne désigne pas nécessairement un token unique.
Il peut représenter un segment relationnel ou une trajectoire
intermédiaire du maillage global.

\begin{remarque}[Finie description et recherche]
Pour une recherche exécutable, les valeurs candidates des \(c_i\)
doivent avoir une représentation finie et un coût calculable.
Cela n'impose pas que la théorie se limite aux chemins courts :
la borne de longueur ou de coût est un paramètre d'exécution
déclaré, et non une modification de la définition de l'alignement.

Une recherche interrompue par son budget doit être signalée comme
incomplète ; elle ne peut pas être déclarée exhaustive.
\end{remarque}

\section{Maillage multidirectionnel du prompt}

\begin{definition}[Maillage contextuel]
Le maillage contextuel du prompt est un graphe orienté typé :

\[
C_P=
(V_P^C,E_P^C,\tau_P^C,\ell_P^C).
\]

Il contient au minimum :
\begin{itemize}
  \item les occurrences \(o_i\) des tokens du prompt ;
  \item les variables contextuelles \(c_i\) ;
  \item les relations d'ordre ou d'ancrage ;
  \item les relations contextuelles multidirectionnelles ;
  \item les rôles nécessaires à la construction de l'alignement.
\end{itemize}
\end{definition}

La multidirectionnalité signifie que les relations contextuelles
ne sont pas réduites à une unique chaîne de voisinage. Elle autorise
plusieurs orientations relationnelles entre les éléments du prompt.

Elle ne signifie pas que toutes les arêtes sont causalement
équivalentes. Chaque arête conserve son type.

\section{Schéma général}

Pour \(P=(x_1,\ldots,x_n)\), un schéma d'ancrage possible est :

\[
c_0\to x_1\to c_1\to x_2
\to\cdots\to c_{n-1}\to x_n\to c_n.
\]

Les \(c_i\) représentent les emplacements contextuels susceptibles
d'accepter des tokens, ensembles de tokens ou trajectoires
intermédiaires lors de la recherche dans \(G_{\Dataset}\).

Le schéma d'ancrage n'est pas, à lui seul, le maillage complet.
Les relations multidirectionnelles supplémentaires appartiennent
à \(E_P^C\).

\section{Exemple : \texttt{je vais bien}}

Soit :

\[
P=(\texttt{je},\texttt{vais},\texttt{bien}).
\]

Le schéma d'ancrage est :

\[
c_1\to\texttt{je}\to c_2
\to\texttt{vais}\to c_3
\to\texttt{bien}\to c_4.
\]

Une partie des relations multidirectionnelles peut être représentée
par :

\[
\begin{aligned}
c_0&\to\texttt{je},\\
c_0&\to\texttt{vais},\\
c_0&\to\texttt{bien},\\
c_0&\to c_4,
\end{aligned}
\]

puis :

\[
\begin{aligned}
\texttt{je}&\to c_1,\\
\texttt{je}&\to\texttt{vais},\\
\texttt{je}&\to\texttt{bien},\\
\texttt{je}&\to c_4,
\end{aligned}
\]

\[
\begin{aligned}
\texttt{vais}&\to c_2,\\
\texttt{vais}&\to\texttt{je},\\
\texttt{vais}&\to\texttt{bien},\\
\texttt{vais}&\to c_4,
\end{aligned}
\]

et :

\[
\begin{aligned}
\texttt{bien}&\to c_3,\\
\texttt{bien}&\to\texttt{vais},\\
\texttt{bien}&\to\texttt{je},\\
\texttt{bien}&\to c_4.
\end{aligned}
\]

Ces relations expriment la construction contextuelle décrite
dans la présente théorie. Leurs types exacts doivent être enregistrés
dans \(\ell_P^C\) ; le simple symbole \(\to\) ne suffit pas à décider
si une relation est contextuelle, ordinale ou causale.

\begin{remarque}[Convention d'indexation]
Le schéma ci-dessus utilise les indices \(c_i\) pour les positions
contextuelles et \(c_0\) pour un emplacement de départ supplémentaire
dans l'exemple multidirectionnel. Dans une implémentation, chaque
rôle doit être identifié par un identifiant typé, afin qu'un même
symbole graphique ne provoque pas de collision entre deux fonctions.
\end{remarque}

\section{Ordre causal du processus de réponse}

La trajectoire :

\[
c_0\leadsto c_x
\]

désigne l'orientation causale globale du processus de recherche et
de construction de la réponse.

Le symbole \(\leadsto\) est réservé à cette fonction d'orientation
causale du processus. Il ne doit pas être confondu avec toute arête
contextuelle de \(C_P\), ni avec toute arête causale d'un écosystème
global.

\begin{definition}[Trajectoire causale de réponse]
Une trajectoire causale de réponse est une trajectoire finie
\(\rho\) dans le maillage global, munie d'un point de départ
contextuellement ancré et d'un point d'aboutissement associé à
une sortie candidate.

On écrit :

\[
\rho:c_0\leadsto c_x.
\]

La trajectoire doit conserver son ordre, les types des arêtes
parcourues et les correspondances avec les ancrages du prompt.
\end{definition}

\begin{remarque}
L'existence d'une trajectoire orientée dans un graphe n'est pas
automatiquement une preuve de causalité physique. Dans cette
extension, elle représente d'abord l'ordre causal interne défini
par le modèle. Toute interprétation plus forte exige des conditions
supplémentaires.
\end{remarque}


% ============================================================
\chapter{Correspondance entre cerveau contextuel et maillage global}
% ============================================================

\section{Fonction de correspondance lexicale}

\begin{definition}[Correspondance lexicale]
Soit \(C_P\) le cerveau contextuel du prompt et \(G\) le maillage
global. La correspondance lexicale est une relation partielle :

\[
M_P\subseteq O_P\times V_G.
\]

Elle relie une occurrence du prompt à un nœud lexical global lorsque
les deux partagent la même identité canonique.

Si l'identité canonique est unique, on peut écrire :

\[
M_P(o_i,v)
\quad\Longleftrightarrow\quad
\pi_{\mathrm{tok}}(o_i)=v.
\]
\end{definition}

Cette correspondance n'intègre pas \(C_P\) dans \(G\). Elle établit
une interface de recherche entre les deux structures.

\section{Correspondance des variables contextuelles}

\begin{definition}[Instanciation contextuelle]
Une instanciation des variables contextuelles est une application
partielle :

\[
\eta:\{c_i\}\rightharpoonup\mathcal{A}(G).
\]

Chaque valeur \(\eta(c_i)\), lorsqu'elle est définie, est une
structure finiment représentée du maillage global.

L'instanciation est admissible si elle respecte les types de valeurs,
les contraintes de chemin, les limites de ressources déclarées
et les règles de compatibilité du protocole.
\end{definition}

L'instanciation permet à un emplacement contextuel d'accepter
un token, un ensemble de tokens ou une trajectoire intermédiaire.

\section{Alignement contextuel}

\begin{definition}[Alignement contextuel d'une trajectoire]
Soit le prompt tokenisé \(P=(x_1,\ldots,x_n)\) et une trajectoire
candidate \(\rho\) du maillage global.

Un alignement contextuel de \(P\) sur \(\rho\) est une application
\[
\alpha:\{1,\ldots,n\}\to\{0,\ldots,|\rho|\}
\]
telle que :

\begin{enumerate}
  \item chaque position \(i\) est associée à une occurrence ou
        un nœud global correspondant à \(x_i\) ;
  \item les positions alignées respectent l'ordre déclaré du prompt
        lorsque l'ordre est une contrainte du protocole ;
  \item entre deux positions alignées, la trajectoire peut contenir
        des tokens ou sous-trajectoires intermédiaires ;
  \item les segments intermédiaires satisfont les contraintes
        contextuelles et relationnelles applicables ;
  \item chaque association est traçable jusqu'aux nœuds et arêtes
        effectivement présents dans le maillage global.
\end{enumerate}
\end{definition}

L'alignement n'exige donc pas que les tokens du prompt soient
directement voisins dans le maillage global. Il exige qu'ils puissent
être retrouvés dans une structure compatible, avec les éléments
intermédiaires éventuellement nécessaires.

\begin{remarque}[Ordre des ancrages]
L'ordre des ancrages est une contrainte par défaut pour une
trajectoire qui représente la progression du prompt. Si une tâche
particulière exige un autre ordre, cette règle doit être explicitement
modifiée dans le protocole. La multidirectionnalité contextuelle
ne doit pas effacer silencieusement l'ordre des ancrages.
\end{remarque}

\section{Alignement sans correspondance directe}

\begin{exemple}
Pour le prompt abstrait :

\[
P=(a,b,c),
\]

une trajectoire globale peut avoir la forme :

\[
c_0\leadsto a\leadsto
(u_1,u_2)\leadsto b\leadsto
(v_1,v_2,v_3)\leadsto c\leadsto c_x.
\]

Les ensembles intermédiaires \((u_1,u_2)\) et
\((v_1,v_2,v_3)\) représentent des tokens ou des segments
relationnels admissibles.

L'alignement est possible si les trois ancrages \(a,b,c\) sont
retrouvés et si les segments intermédiaires respectent les contraintes
du maillage contextuel et de l'écosystème global.
\end{exemple}

\begin{remarque}
La présence des tokens \(a,b,c\) quelque part dans le maillage
global n'est pas suffisante. Ils doivent être retrouvés dans une
configuration relationnelle compatible, et non simplement dans
trois endroits sans relation admissible.
\end{remarque}


% ============================================================
\part{Recherche d'une réponse}
% ============================================================

\chapter{Architecture de recherche et conditions de cohérence}

\section{Problème de recherche}

Soient :

\begin{itemize}
  \item \(P\), le prompt courant ;
  \item \(C_P\), son maillage contextuel ;
  \item \(G\), le maillage global ;
  \item \(\mathcal{A}(G)\), l'ensemble des structures globales
        admissibles pour les variables contextuelles ;
  \item \(\mathcal{Y}\), l'espace des réponses candidates.
\end{itemize}

Le système doit chercher une trajectoire globale qui permet de
retrouver les ancrages du prompt, d'exploiter les écosystèmes des
tokens et d'aboutir à une réponse candidate selon l'orientation
causale définie.

\begin{definition}[Trajectoire candidate de réponse]
Une trajectoire candidate est un tuple :

\[
\rho=(v_0,e_1,v_1,\ldots,e_m,v_m),
\]

où les \(v_i\) sont des nœuds du maillage global et les \(e_i\)
des arêtes orientées reliant ces nœuds.

La trajectoire est admissible si :
\begin{enumerate}
  \item ses arêtes existent dans \(G\) ;
  \item ses types relationnels sont conservés ;
  \item ses contraintes bloquantes sont satisfaites ;
  \item ses ancrages contextuels sont définis ;
  \item son coût est fini et calculable ;
  \item sa provenance est disponible.
\end{enumerate}
\end{definition}

\section{Les trois niveaux de contraintes}

\subsection{Niveau I : alignement contextuel}

Le premier niveau exige que les éléments du prompt puissent être
alignés avec une trajectoire globale.

On note :

\[
\mathsf{A}_P(\rho)=1
\]

si un alignement contextuel admissible existe, et \(0\) sinon.

L'absence d'alignement signifie que la trajectoire n'a pas établi
la correspondance structurelle requise avec le prompt.

\subsection{Niveau II : écosystème causal et anti-causal}

Le deuxième niveau examine les relations globales parcourues par
la trajectoire, y compris les relations incidentes aux tokens
alignés et les sous-trajectoires intermédiaires.

On définit :

\[
\mathsf{E}_G(\rho)\in
\{\mathsf{compatible},
\mathsf{incompatible},
\mathsf{indéterminé}\}.
\]

Le statut est compatible si toutes les contraintes relationnelles
bloquantes applicables sont satisfaites et si les contraintes
relationnelles requises sont présentes.

Il est incompatible lorsqu'une contrainte bloquante certifiée est
violée.

Il est indéterminé lorsque les informations nécessaires ne suffisent
pas à établir l'un des deux statuts précédents.

\subsection{Niveau III : orientation causale}

Le troisième niveau vérifie que la trajectoire candidate appartient
à une progression orientée compatible avec le processus de réponse :

\[
\mathsf{O}(\rho)=1.
\]

Ce statut exige que la trajectoire possède un point de départ
contextuellement ancré, un ordre de progression défini et un point
d'aboutissement associé à la réponse candidate.

L'orientation n'est pas déduite de la seule présence des bons tokens.

\section{Admissibilité structurelle}

\begin{definition}[Trajectoire structurellement admissible]
Une trajectoire \(\rho\) est structurellement admissible pour le
prompt \(P\) lorsque :

\[
\mathsf{A}_P(\rho)=1,
\]

\[
\mathsf{E}_G(\rho)=\mathsf{compatible},
\]

et

\[
\mathsf{O}(\rho)=1.
\]

L'ensemble des trajectoires admissibles est :

\[
\mathcal{T}(P,G)
=
\left\{
\rho:
\mathsf{A}_P(\rho)=1
\land
\mathsf{E}_G(\rho)=\mathsf{compatible}
\land
\mathsf{O}(\rho)=1
\right\}.
\]
\end{definition}

\begin{remarque}
Une trajectoire déclarée indéterminée ne doit pas être confondue
avec une trajectoire prouvée incompatible. L'absence de preuve
d'admissibilité n'est pas toujours une preuve d'inadmissibilité.
Cette distinction est nécessaire pour conserver un statut correct
aux recherches interrompues ou insuffisamment documentées.
\end{remarque}

\section{Pourquoi l'alignement seul ne suffit pas}

\begin{proposition}[Insuffisance de l'alignement isolé]
L'existence d'un alignement contextuel ne garantit pas, à elle seule,
l'admissibilité structurelle d'une trajectoire.

\end{proposition}

\begin{proof}
Par définition, l'alignement ne vérifie que la correspondance
entre les ancrages du prompt et la trajectoire candidate. Il ne
garantit ni la compatibilité des relations causales et anti-causales,
ni l'orientation causale globale de la trajectoire. Une trajectoire
peut donc satisfaire \(\mathsf{A}_P(\rho)=1\) tout en ayant
\(\mathsf{E}_G(\rho)=\mathsf{incompatible}\), ou
\(\mathsf{O}(\rho)=0\). Elle n'appartient alors pas à
\(\mathcal{T}(P,G)\).
\end{proof}

\section{Rôle du maillage global dans la discrimination}

Le maillage global n'est pas un ensemble non structuré de tokens.
Chaque token fusionné peut être associé à un écosystème de relations
typées, avec leur provenance et leurs statuts.

La recherche d'une réponse doit donc tenir compte simultanément :

\begin{enumerate}
  \item de l'identité des tokens retrouvés ;
  \item des chemins qui les relient ;
  \item des relations naturelles et anti-causales parcourues ;
  \item des relations contextuelles imposées par le prompt ;
  \item de l'orientation causale du parcours complet.
\end{enumerate}

Une réponse arbitraire qui n'a aucun alignement ni aucune trajectoire
compatible ne peut pas être retenue comme réponse admissible par
ce mécanisme.

Cette propriété est structurelle. Elle ne signifie pas que toutes
les réponses admissibles sont nécessairement vraies, ni qu'une
réponse fausse ne peut jamais être compatible avec un maillage
inexact ou incomplet.

\section{La réponse n'a pas à répéter les tokens du prompt}

\begin{remarque}
L'alignement contextuel porte sur la relation entre le prompt et la
trajectoire de recherche. Il ne requiert pas que la réponse textuelle
reproduise les tokens du prompt.

Pour une question \(P=(a,b,c)\), une réponse \(R=(d,e,f)\) peut être
admissible si une trajectoire globale relie la structure alignée sur
\(a,b,c\) à la structure générative correspondant à \(d,e,f\), et si
les contraintes contextuelles, relationnelles et causales sont
satisfaites.

À l'inverse, le seul fait que \(d,e,f\) soient des tokens connus du
maillage global ne suffit pas à justifier leur émission comme réponse.
\end{remarque}


% ============================================================
\chapter{Règles de sélection et construction de la sortie}
% ============================================================

\section{Ensemble des réponses candidates}

Soit \(\mathcal{T}(P,G)\) l'ensemble des trajectoires structurellement
admissibles.

Une application de lecture de sortie :

\[
\operatorname{out}:\mathcal{T}(P,G)\to\mathcal{Y}
\]

associe à chaque trajectoire une réponse candidate.

L'ensemble des réponses candidates structurellement admissibles est :

\[
\Cand(P,G)
=
\{
\operatorname{out}(\rho):
\rho\in\mathcal{T}(P,G)
\}.
\]

Cette définition sépare la trajectoire interne et le texte produit.
Une trajectoire est un objet du maillage ; une réponse est sa
réalisation textuelle selon une règle de génération déclarée.

\section{Discrimination par l'écosystème}

\begin{definition}[Signature relationnelle d'une trajectoire]
La signature relationnelle d'une trajectoire \(\rho\) est la structure
qui conserve, dans l'ordre du parcours :

\begin{itemize}
  \item les identités des nœuds traversés ;
  \item les types des arêtes ;
  \item les relations causales naturelles ;
  \item les relations anti-causales ;
  \item les contraintes contextuelles satisfaites ;
  \item les provenances disponibles ;
  \item les éventuelles violations ou incertitudes.
\end{itemize}
\end{definition}

Deux trajectoires ayant les mêmes tokens visibles peuvent avoir des
signatures relationnelles différentes. Le système ne doit donc pas
les considérer comme équivalentes au seul motif que leurs sorties
textuelles se ressemblent.

\section{Règle de sélection}

\begin{postulat}[Sélection interne des trajectoires]
Le système sélectionne d'abord parmi les trajectoires qui satisfont
les contraintes bloquantes. Il peut ensuite ordonner les trajectoires
admissibles selon les relations contextuelles, les écosystèmes
causaux et anti-causaux, la provenance et les coûts déclarés.

Le mécanisme de classement doit conserver la distinction entre :
\begin{enumerate}
  \item les contraintes impératives ;
  \item les indices relationnels préférentiels ;
  \item les incertitudes non résolues.
\end{enumerate}
\end{postulat}

Pour un classement numérique, on peut définir une fonction :

\[
\Score(\rho\mid P,G)
=
F\bigl(
A(\rho),
E_{\Nat}(\rho),
E_{\Anti}(\rho),
O(\rho),
\operatorname{Prov}(\rho),
\cost(\rho)
\bigr).
\]

La fonction \(F\) doit être explicitement définie dans toute
implémentation. La présente théorie n'impose pas arbitrairement
des coefficients numériques qui ne seraient pas justifiés par le
protocole.

\begin{remarque}[Critère interne ou critère extérieur]
Le classement peut exploiter les structures causales déjà présentes
dans le maillage global. Il n'est donc pas nécessaire de supposer
qu'un critère de pertinence entièrement indépendant doit être ajouté
dans tous les cas.

En revanche, si plusieurs trajectoires satisfont exactement les
contraintes et possèdent des signatures équivalentes pour toutes les
observations utilisées, le mécanisme déclaré ne peut pas les
départager par une information qu'il ne possède pas. Une règle
déterministe de départage peut choisir une représentante, mais ce
choix ne constitue pas une preuve de supériorité sémantique.
\end{remarque}

\section{Sélection en cas de pluralité}

Lorsque plusieurs trajectoires sont admissibles, le protocole peut
appliquer, dans l'ordre :

\begin{enumerate}
  \item l'élimination des violations de contraintes bloquantes ;
  \item l'évaluation de la compatibilité des écosystèmes ;
  \item la comparaison des alignements contextuels ;
  \item la comparaison des signatures relationnelles ;
  \item la comparaison des provenances et statuts de validation ;
  \item un classement selon la fonction de score déclarée ;
  \item une règle de départage déterministe, si elle est nécessaire.
\end{enumerate}

La dernière règle ne doit pas être présentée comme une preuve
de vérité. Elle ne fait que sélectionner une sortie parmi les
candidates qui restent équivalentes selon le protocole.

\section{Statut d'absence de réponse}

\begin{definition}[Statuts de recherche]
Une recherche doit distinguer au minimum :

\begin{description}[leftmargin=3.2cm,style=nextline]
  \item[FOUND] Une trajectoire admissible et sa sortie ont été
        trouvées.
  \item[NO\_MATCH] Aucune correspondance n'existe dans le domaine
        fini effectivement exhausté.
  \item[INCOMPATIBLE] Les trajectoires considérées violent des
        contraintes bloquantes certifiées.
  \item[UNKNOWN] Les données disponibles ne permettent pas de
        décider l'admissibilité.
  \item[INCOMPLETE] La recherche s'est arrêtée avant l'épuisement
        de son domaine déclaré.
  \item[RESOURCE\_STOP] L'arrêt est dû au budget de calcul ou
        à une limite d'exécution.
\end{description}
\end{definition}

\begin{remarque}
Le statut \texttt{NO\_MATCH} exige l'exhaustion du domaine déclaré.
Une recherche interrompue ou gloutonne qui n'a pas trouvé de réponse
ne prouve pas l'inexistence d'une réponse.
\end{remarque}


% ============================================================
\part{Propriétés formelles et limites logiques}
% ============================================================

\chapter{Propriétés démontrables de l'architecture}

\section{Séparation des structures}

\begin{proposition}[Non-modification du maillage global par requête]
Soit \(G\) le maillage global avant une requête \(P\), et \(C_P\)
le cerveau contextuel construit pour cette requête.

Si l'opération d'inférence n'effectue que des lectures de \(G\) et
construit \(C_P\) séparément, alors :

\[
G_{\mathrm{après}}=G_{\mathrm{avant}}.
\]

\end{proposition}

\begin{proof}
Par hypothèse, l'opération d'inférence ne possède aucune opération
d'écriture sur \(G\). La construction de \(C_P\) agit sur une
structure distincte. Aucune modification des composantes de \(G\)
n'est donc produite par cette opération.
\end{proof}

\section{Monotonie de l'ensemble des contraintes bloquantes}

\begin{proposition}[Réduction de l'ensemble admissible]
Soit \(T\) un ensemble fini de trajectoires candidates. Soit
\(\mathcal{K}_1\) un ensemble de contraintes bloquantes, puis
\(\mathcal{K}_2\) un ensemble qui les contient :

\[
\mathcal{K}_1\subseteq\mathcal{K}_2.
\]

Si une trajectoire est admissible relativement à un ensemble de
contraintes lorsque toutes ces contraintes sont satisfaites, alors :

\[
T_{\mathcal{K}_2}\subseteq T_{\mathcal{K}_1}.
\]

\end{proposition}

\begin{proof}
Toute trajectoire satisfaisant toutes les contraintes de
\(\mathcal{K}_2\) satisfait celles de \(\mathcal{K}_1\), puisque
ces dernières sont incluses dans \(\mathcal{K}_2\).
L'inclusion des ensembles admissibles en découle.
\end{proof}

Cette propriété établit qu'une contrainte bloquante supplémentaire
ne peut pas élargir l'ensemble des trajectoires admissibles. Elle
ne prouve pas que toute contrainte ajoutée est correcte : la
validité de chaque contrainte doit être établie séparément.

\section{Propriété de rejet d'une réponse sans alignement}

\begin{proposition}[Rejet des sorties sans trajectoire admissible]
Si une réponse \(R\) n'est l'image d'aucune trajectoire de
\(\mathcal{T}(P,G)\), alors :

\[
R\notin\Cand(P,G).
\]

\end{proposition}

\begin{proof}
Par définition,

\[
\Cand(P,G)
=
\{\operatorname{out}(\rho):
\rho\in\mathcal{T}(P,G)\}.
\]

Si aucune trajectoire admissible ne produit \(R\), alors \(R\)
n'appartient pas à cet ensemble.
\end{proof}

\begin{remarque}
Cette proposition est une conséquence de la définition du mécanisme
de sélection. Elle ne démontre pas qu'une implémentation concrète
respecte automatiquement la règle ; cela doit être vérifié par
l'exécution, les tests et la traçabilité.
\end{remarque}

\section{Rôle des relations causales dans la cohérence}

\begin{proposition}[Effet des contraintes causales certifiées]
Supposons que certaines trajectoires soient incompatibles avec
des contraintes causales ou anti-causales certifiées, et que le
système applique effectivement ces contraintes avant la sélection.

Alors aucune de ces trajectoires incompatibles ne peut appartenir
à l'ensemble admissible.

\end{proposition}

\begin{proof}
L'admissibilité exige la satisfaction de toutes les contraintes
bloquantes. Toute trajectoire qui en viole au moins une est exclue
par définition.
\end{proof}

\begin{remarque}[Portée du résultat]
Le résultat démontre un effet de filtrage structurel, et non une
garantie absolue de vérité. Il dépend de la correction des contraintes
et de leur exécution effective.

Si le maillage global contient une relation erronée, si une
annotation est inadéquate ou si une relation anti-causale est
incorrectement déclarée bloquante, une réponse correcte pourrait
être écartée ou une réponse incorrecte conservée.
\end{remarque}

\section{Équivalence observationnelle des trajectoires}

\begin{definition}[Indiscernabilité de deux trajectoires]
Soit \(\mathcal{Q}\) une famille de sondes appliquées aux trajectoires.
Deux trajectoires \(\rho_1,\rho_2\) sont observationnellement
indiscernables relativement à \(\mathcal{Q}\) lorsque :

\[
\rho_1\ObsRel_{\mathcal{Q}}\rho_2
\quad\Longleftrightarrow\quad
I_{\mathcal{Q}}(\rho_1)=I_{\mathcal{Q}}(\rho_2).
\]
\end{definition}

\begin{proposition}[Limite de discrimination observationnelle]
Si deux trajectoires sont indiscernables pour toutes les sondes
effectivement utilisées par le mécanisme de sélection, toute règle
déterministe qui dépend uniquement de ces empreintes leur attribue
la même valeur de classement.

\end{proposition}

\begin{proof}
Les entrées de la règle de classement sont identiques pour les deux
trajectoires. Le déterminisme impose donc des sorties identiques.
\end{proof}

Cette proposition est compatible avec le rôle de l'écosystème
global : si l'écosystème apporte des sondes relationnelles qui
distinguent les trajectoires, elles peuvent raffiner l'empreinte.
Si aucune observation disponible ne les distingue, la théorie
ne permet pas d'inventer cette distinction.

\section{Cohérence structurelle et vérité factuelle}

\begin{definition}[Cohérence structurelle]
Une réponse est structurellement cohérente relativement à \(P\),
\(G\) et au protocole \(\Pi\) si elle est produite par une trajectoire
admissible selon les contraintes d'alignement, de relations et
d'orientation de \(\Pi\).
\end{definition}

\begin{definition}[Vérité factuelle relativement à une tâche]
Une réponse est factuellement correcte relativement à une tâche
lorsqu'elle satisfait le critère de vérité ou de correction externe
défini pour cette tâche, dans le domaine concerné.
\end{definition}

\begin{proposition}[Non-équivalence logique]
La cohérence structurelle ne suffit pas, sans hypothèse supplémentaire,
à démontrer la vérité factuelle d'une réponse.
\end{proposition}

\begin{proof}
La cohérence structurelle dépend de la structure du maillage et
des règles du protocole. La vérité factuelle dépend d'un critère
de correction relatif à la tâche. En l'absence d'une hypothèse
établissant que le maillage et ses relations représentent
correctement toutes les distinctions factuelles nécessaires,
la première propriété n'implique pas logiquement la seconde.
\end{proof}

\begin{remarque}
Cette distinction ne réduit pas le rôle du maillage global. Elle
délimite précisément ce que l'architecture peut garantir par sa
structure et ce qui nécessite une validation supplémentaire des
données et des relations qui la constituent.
\end{remarque}


% ============================================================
\chapter{Algorithme de construction et de recherche}
% ============================================================

\section{Construction du maillage global}

Le pseudo-algorithme suivant décrit le contrat logique de construction.
Les opérations internes de tokenisation, d'annotation et de validation
doivent être définies par l'implémentation utilisée.

\begin{verbatim}
CONSTRUIRE_MAILLAGE_GLOBAL(dataset, protocole):

    initialiser G avec des tables vides
    initialiser le registre de provenance

    pour chaque exemple z du dataset:

        vérifier que z possède un identifiant stable
        enregistrer sa source et son statut

        tokeniser le prompt P_z
        tokeniser la réponse R_z

        créer des occurrences distinctes
        pour chaque position du prompt et de la réponse

        associer chaque occurrence à son identité lexicale
        selon la normalisation déclarée

        construire le cerveau local Brain_z

        construire les relations observables et contextuelles
        selon les règles déclarées

        attribuer un type causal, anti-causal ou inconnu
        uniquement si le protocole le permet

        conserver la provenance de chaque relation

        appliquer la politique d'admission du dataset

        si les contributions sont admissibles:
            fusionner les identités lexicales communes dans G
            agréger les relations en conservant :
                leur type
                leur direction
                leur poids
                leur provenance
                leur statut de validation

        conserver Brain_z comme structure reconstructible

    retourner G et le registre de provenance
\end{verbatim}

\section{Construction du cerveau contextuel}

\begin{verbatim}
CONSTRUIRE_CERVEAU_CONTEXTUEL(prompt P, protocole):

    tokeniser P en occurrences distinctes
    conserver l'identité lexicale de chaque occurrence

    créer les nœuds d'occurrence
    créer les variables contextuelles c_i

    construire les relations d'ordre et d'ancrage
    construire les relations contextuelles multidirectionnelles
    typer les relations selon le protocole

    conserver les occurrences, les positions et les rôles
    ne pas modifier le maillage global

    retourner C_P
\end{verbatim}

\section{Recherche d'une réponse}

\begin{verbatim}
RECHERCHER_REPONSE(P, G, protocole, budget):

    C_P <- CONSTRUIRE_CERVEAU_CONTEXTUEL(P, protocole)

    initialiser la frontière de recherche
    initialiser les trajectoires candidates
    initialiser les statuts de provenance et d'incertitude

    rechercher les correspondances lexicales
    rechercher les alignements contextuels possibles

    pour chaque trajectoire candidate rho:

        vérifier l'existence des arêtes globales
        vérifier l'alignement des tokens du prompt
        autoriser les segments intermédiaires admissibles

        examiner les écosystèmes des tokens alignés
        examiner les relations causales naturelles
        examiner les relations anti-causales

        appliquer les contraintes anti-causales bloquantes
        réorienter l'exploration selon les contraintes souples

        vérifier l'orientation globale c_0 -> c_x
        vérifier la provenance et les statuts de validation

        si une contrainte bloquante est violée:
            rejeter rho
        sinon si les informations sont insuffisantes:
            marquer rho comme INDETERMINE
        sinon si toutes les conditions sont satisfaites:
            ajouter rho aux trajectoires admissibles

        arrêter uniquement selon le budget déclaré
        et conserver le statut exact de l'arrêt

    si une trajectoire admissible est trouvée:
        sélectionner selon le protocole
        produire la réponse par la règle de sortie déclarée
        retourner FOUND avec la trajectoire et sa provenance

    sinon si la frontière a été exhaustée:
        retourner NO_MATCH ou INCOMPATIBLE selon le cas

    sinon:
        retourner INCOMPLETE ou RESOURCE_STOP
\end{verbatim}

\section{Contrat d'implémentation}

Une implémentation conforme doit permettre de répondre aux questions
suivantes pour toute réponse produite :

\begin{enumerate}
  \item Quel prompt a généré la recherche ?
  \item Quel cerveau contextuel a été construit ?
  \item Quels tokens ont été alignés ?
  \item Quelles variables contextuelles ont été instanciées ?
  \item Quelle trajectoire globale a été retenue ?
  \item Quelles arêtes causales et anti-causales ont été parcourues ?
  \item Quelles contraintes ont été satisfaites ou violées ?
  \item Quelles sources ont contribué aux relations utilisées ?
  \item Quel critère a départagé les trajectoires candidates ?
  \item La recherche a-t-elle été exhaustive dans son domaine déclaré ?
\end{enumerate}

Si une réponse à l'une de ces questions est requise pour justifier
la sortie mais n'est pas disponible, la traçabilité correspondante
doit être marquée comme manquante plutôt que reconstruite
rétroactivement par conjecture.


% ============================================================
\part{Résolution observationnelle et validation}
% ============================================================

\chapter{Résolution, seuils et raffinement des cerveaux}

\section{Empreinte d'une structure relationnelle}

Le cadre observationnel hérité peut être appliqué aux cerveaux et
aux trajectoires.

Soit \(\mathcal{X}_G\) un ensemble de structures candidates issues
du maillage global. Une famille de sondes relationnelles
\(\mathcal{F}_\lambda\) peut observer, par exemple :

\begin{itemize}
  \item l'identité des tokens ;
  \item les types d'arêtes ;
  \item les relations entre occurrences ;
  \item la présence de contraintes causales ;
  \item la provenance des relations ;
  \item la compatibilité avec le cerveau contextuel ;
  \item la structure de la trajectoire de réponse.
\end{itemize}

Pour une structure \(x\in\mathcal{X}_G\), son empreinte à l'échelle
\(\lambda\) est :

\[
I_{\Phi,\lambda}^{P}(x)
=
\bigl(
\Phi(P,x,f)
\bigr)_{f\in\mathcal{F}_\lambda}.
\]

\section{Raffinement observationnel}

Si :

\[
\mathcal{F}_{\lambda}
\subseteq
\mathcal{F}_{\mu},
\]

alors l'empreinte à l'échelle \(\mu\) contient toutes les
observations de l'échelle \(\lambda\), ainsi que les observations
supplémentaires.

\begin{proposition}[Raffinement des candidats indiscernables]
Si l'empreinte à l'échelle \(\mu\) est un enrichissement de celle
de \(\lambda\), alors deux structures indiscernables à l'échelle
\(\mu\) sont indiscernables à l'échelle \(\lambda\), mais la
réciproque n'est pas garantie.
\end{proposition}

\begin{proof}
L'empreinte enrichie contient toutes les coordonnées de l'empreinte
initiale. L'égalité de toutes les coordonnées enrichies implique
donc l'égalité des coordonnées initiales.
\end{proof}

Cette propriété donne une interprétation observationnelle au
raffinement du cerveau contextuel et des sondes de trajectoire :
des relations supplémentaires peuvent distinguer des candidats
qui restaient indiscernables avec les observations précédentes.

\section{Seuils de résolution}

Lorsque les sondes sont organisées en une famille croissante,
les seuils de résolution du socle peuvent être utilisés pour
représenter les niveaux auxquels des candidats deviennent
discernables.

La présente extension ne postule pas une métrique lexicale
universelle. Les seuils doivent être construits à partir d'une
fonction de séparation déclarée et de ses hypothèses.

Pour deux candidats \(x,y\), on peut définir une fonction de seuil
\(K_P(x,y)\) relativement au protocole d'observation choisi. Sa
signification dépend de la famille de sondes, de l'échelle et de
la convention de bord déclarées.

\begin{remarque}
Une distance entre tokens ne doit pas être introduite simplement
par analogie avec une distance géométrique. Si une distance,
un coût ou un seuil est utilisé dans la recherche, sa construction
doit être spécifiée séparément.
\end{remarque}


% ============================================================
\chapter{Validation empirique et tests de falsification}
% ============================================================

\section{Objectif de la validation}

La validation doit tester si l'architecture améliore effectivement
la pertinence des réponses et réduit les sorties sans alignement,
dans les conditions expérimentales déclarées.

Il faut distinguer :

\begin{enumerate}
  \item la validité des définitions ;
  \item la correction de l'implémentation ;
  \item la qualité du dataset ;
  \item la fidélité des annotations causales et anti-causales ;
  \item l'efficacité de la recherche ;
  \item la qualité sémantique et factuelle des réponses.
\end{enumerate}

Un résultat sur l'une de ces dimensions ne constitue pas
automatiquement une preuve sur les autres.

\section{Ablations minimales}

Les tests doivent comparer au moins les variantes suivantes :

\begin{description}[leftmargin=3.2cm,style=nextline]
  \item[Sans maillage contextuel]
        Recherche sans structure multidirectionnelle du prompt.

  \item[Sans fusion lexicale]
        Conservation de nœuds lexicaux séparés pour les occurrences,
        afin de mesurer l'effet de la fusion des identités.

  \item[Sans écosystème]
        Recherche sans exploitation des relations incidentes
        aux tokens.

  \item[Sans anti-causalité]
        Désactivation des contraintes et mécanismes anti-causaux.

  \item[Sans ordre causal global]
        Retrait de l'orientation \(c_0\leadsto c_x\).

  \item[Sans segments intermédiaires]
        Alignement direct uniquement, afin de mesurer l'apport
        des trajectoires intermédiaires.

  \item[Sans provenance]
        Variante de diagnostic seulement : elle mesure le coût
        de la perte de traçabilité, mais ne constitue pas une
        configuration de référence recommandée.
\end{description}

\section{Mesures à enregistrer}

Pour chaque variante, il faut enregistrer :

\begin{itemize}
  \item le taux de réponses disposant d'un alignement traçable ;
  \item le taux de violations des contraintes bloquantes ;
  \item le taux de réponses jugées pertinentes selon un protocole
        d'évaluation défini indépendamment ;
  \item le taux de réponses factuellement correctes lorsque cette
        propriété est évaluable ;
  \item le nombre de trajectoires explorées ;
  \item le coût de recherche et le nombre de relations parcourues ;
  \item la proportion de recherches terminées, incomplètes ou arrêtées
        pour ressources ;
  \item la fréquence des conflits entre annotations causales
        et anti-causales ;
  \item la stabilité des résultats sur plusieurs versions du dataset.
\end{itemize}

Les évaluations de pertinence et de vérité doivent préciser leur
protocole et leurs critères. Elles ne doivent pas être remplacées
par le seul score interne du maillage.

\section{Tests adversariaux}

Les cas suivants doivent être inclus :

\begin{enumerate}
  \item des prompts comportant des négations ;
  \item des prompts comportant des tokens ambigus ;
  \item des prompts avec répétitions de tokens ;
  \item des prompts dont les tokens existent dans le maillage
        mais sans trajectoire compatible ;
  \item des prompts dont la réponse attendue utilise des tokens
        absents du prompt ;
  \item des cas où plusieurs trajectoires sont admissibles ;
  \item des cas où une trajectoire semble plausible mais viole
        une contrainte anti-causale certifiée ;
  \item des cas où les relations du dataset sont contradictoires ;
  \item des cas où aucune réponse admissible n'est retrouvée ;
  \item des cas où la recherche s'arrête avant l'exhaustion.
\end{enumerate}

\section{Critère de réussite}

La réussite de l'architecture ne doit pas être définie par le seul
fait que le système produit une réponse.

Elle exige des résultats observables sur la qualité des alignements,
la conformité aux contraintes, la traçabilité, les coûts de recherche
et la qualité de la sortie, relativement à des références et à un
protocole d'évaluation préalablement définis.


% ============================================================
\chapter{Conditions, limites et statut final}
% ============================================================

\section{Conditions de validité opérationnelle}

L'architecture repose sur les conditions suivantes :

\begin{enumerate}
  \item une tokenisation et une normalisation reproductibles ;
  \item une identité lexicale commune cohérente ;
  \item une construction déterminée des cerveaux locaux ;
  \item une règle d'agrégation qui préserve les types et provenances ;
  \item un protocole d'attribution des relations causales et
        anti-causales ;
  \item une séparation effective entre maillage global et cerveau
        contextuel ;
  \item une définition des valeurs admissibles des variables \(c_i\) ;
  \item des règles explicites d'alignement ;
  \item une orientation causale de recherche définie ;
  \item des contraintes bloquantes et préférentielles typées ;
  \item une règle de sortie de trajectoire ;
  \item des statuts de recherche permettant de distinguer absence
        de réponse, incompatibilité et recherche incomplète.
\end{enumerate}

\section{Limites logiques à ne pas confondre avec des défauts de définition}

\begin{limite}[Données incomplètes]
Un maillage construit à partir d'un dataset fini ne contient pas
nécessairement toutes les relations nécessaires à une réponse.
L'absence de trajectoire retrouvée ne prouve l'absence de toute
trajectoire possible que si le domaine de recherche déclaré a été
exhausté.
\end{limite}

\begin{limite}[Relations incorrectes]
Une relation erronée ou une annotation causale incorrecte peut
orienter la recherche dans une mauvaise direction. La provenance,
la validation et les mécanismes de correction sont donc des
composantes de l'architecture, et non des métadonnées facultatives.
\end{limite}

\begin{limite}[Ambiguïté lexicale]
Une identité lexicale commune n'implique pas une identité de sens
dans tous les contextes. Les occurrences doivent conserver leurs
positions, leurs relations et les contextes d'origine. Si nécessaire,
des sondes contextuelles doivent distinguer les usages.
\end{limite}

\begin{limite}[Trajectoires multiples]
Plusieurs trajectoires peuvent satisfaire toutes les contraintes.
Si le protocole ne distingue pas ces trajectoires, le choix d'une
représentante est une règle de sélection, et non une preuve que
cette trajectoire est la seule correcte.
\end{limite}

\begin{limite}[Vérité factuelle]
La cohérence interne du maillage ne garantit pas que les données
du maillage soient factuellement correctes. La vérité factuelle
exige un critère de validation relatif à la tâche.
\end{limite}

\begin{limite}[Budget de recherche]
Une limite de calcul peut empêcher de trouver une trajectoire
admissible existante. L'arrêt pour ressources ne doit pas être
transformé en preuve d'inexistence ou d'incompatibilité.
\end{limite}

\section{Synthèse formelle}

L'architecture entière peut être résumée par la chaîne de
constructions suivante :

\[
\boxed{
\Dataset
\longrightarrow
\{\Brain_z\}_{z\in\Dataset}
\longrightarrow
G_{\Dataset}
}
\]

où chaque cerveau local est construit à partir d'un exemple
prompt--réponse et où les identités lexicales communes sont
fusionnées en conservant les occurrences, les relations et les
provenances.

Pour un nouveau prompt \(P\), le cerveau contextuel est construit
séparément :

\[
\boxed{
P\longrightarrow C_P
}
\]

Le cerveau contextuel n'est pas fusionné dans le maillage global.
Il définit les alignements et les contraintes de recherche.

La recherche produit des trajectoires candidates :

\[
\boxed{
(C_P,G_{\Dataset})
\longrightarrow
\mathcal{T}(P,G_{\Dataset})
}
\]

où une trajectoire admissible doit satisfaire conjointement :

\[
\boxed{
\begin{aligned}
&\text{alignement contextuel},\\
&\text{compatibilité de l'écosystème global},\\
&\text{contraintes causales et anti-causales},\\
&\text{orientation causale }c_0\leadsto c_x.
\end{aligned}
}
\]

La sortie est ensuite produite par une règle de lecture et de
génération définie :

\[
\boxed{
R=\operatorname{out}(\rho),
\qquad
\rho\in\mathcal{T}(P,G_{\Dataset}).
}
\]

La propriété centrale de cette extension est donc la suivante :

\begin{minipage}{0.94\textwidth}
\emph{
Une réponse n'est pas admissible uniquement parce que ses tokens
existent dans le maillage global. Elle doit être soutenue par une
trajectoire globale alignée sur le contexte du prompt, compatible
avec l'écosystème relationnel des tokens et respectant l'orientation
causale et les contraintes du processus de recherche.
}
\end{minipage}

\section{Statut final de la théorie}

La présente extension définit une architecture relationnelle
cohérente avec le cadre observationnel de la Théorie des Empreintes.
Elle sépare les structures locales du dataset, le maillage global,
le cerveau contextuel d'une requête, les mécanismes d'alignement
et l'orientation causale de la recherche.

Les propriétés directement démontrées dans ce texte sont celles
qui découlent des définitions et des hypothèses explicitement
énoncées : conservation de la provenance sous l'agrégation déclarée,
non-modification du maillage global par une requête en lecture seule,
raffinement des contraintes bloquantes et exclusion des trajectoires
qui violent effectivement ces contraintes.

L'efficacité empirique, la qualité des annotations causales,
la réduction effective des incohérences et la correction factuelle
des réponses sont des propriétés à mesurer et à valider sur des
datasets et des protocoles déclarés. Elles ne sont pas présentées
comme des théorèmes universels en l'absence de ces validations.

\end{document}